\documentclass[10pt,a4paper]{amsart}
\usepackage[margin=19mm]{geometry}
\usepackage{amsmath,amssymb,mathtools}
\usepackage[scr=rsfs,cal=euler]{mathalfa}
\usepackage{xfrac}
\usepackage[unicode,hidelinks,bookmarks=false]{hyperref}
\newcommand{\ie}{i.e.\ }
\newcommand{\cf}{cf.\ }
\newcommand{\resp}{respectively\ }
\newcommand{\Subsection}{\subsection}
\providecommand{\nobreakdash}{\nobreak\hbox{-}}
\setlength{\emergencystretch}{2em}
\multlinegap0pt
\newtheorem{proposition}{Proposition}
\title{Affine Angles via Area Cross Ratio and Isoptic Hyperbolas}
\author{Masanori Nakazato}
\date{Source: arXiv:2603.22218v2 (CC BY 4.0)}
\begin{document}
\maketitle

\noindent\emph{Source and licence.} Affine Angles via Area Cross Ratio and Isoptic Hyperbolas. Version arXiv:2603.22218v2 (CC BY 4.0), distributed by arXiv under the Creative Commons Attribution 4.0 International licence, http://creativecommons.org/licenses/by/4.0/, as recorded in the arXiv licence metadata for this exact version and shown on the abstract page https://arxiv.org/abs/2603.22218v2. Full licence notice: \texttt{licenses/arxiv-2603.22218-CC-BY-4.0.txt}.\par\smallskip

\noindent\textit{Editorial scope.} Counted unit: the article's proposition prop:hyperbola-area-invariance (source0.tex lines 2018--2033). The article's complete original proof of it is delivered with it (source0.tex lines 2035--2155, one \texttt{proof} environment of 121 lines). No mathematics is written, repaired or re-derived: the statement and the proof are the article's own lines, in the article's own notation, and no retained line is substituted or re-flowed.  No further source-local context is needed: every cross-reference of every retained line resolves inside the two delivered spans.  Nothing was omitted from any retained span, so the package carries no omission record.  No retained line contains a citation, so the wrapper carries no bibliography and no bibliography entry.  No retained line contains a figure environment, an image-inclusion command or drawing code, so no image is delivered and the package declares no asset dependency.  Editorial formatting only, in addition: an AMS \texttt{amsart} preamble that restates the article's own declarations and macros used by the retained lines, namely the environments \texttt{proposition} on separate counters, together with the standard packages those lines need.  The retained environments print in this document as Proposition 1 in order of appearance, whereas the article prints the same items with the numbering of its own full document.  The complete native source of the article as distributed in the arXiv e-print for this exact version is bundled unchanged as \texttt{sources/197006/source0.tex}.
\begin{proposition}[Hyperbolic version: area invariance under geometric progression]
\label{prop:hyperbola-area-invariance}
Let $a,r>0$ with $r\neq1$.
On the hyperbola
\[
\mathcal H:\ xy=\kappa \qquad (x>0),
\]
take four distinct points $A,B,C,D$ such that
\[
x_A=a,\quad x_B=ar,\quad x_C=ar^2,\quad x_D=ar^3.
\]
Let $P$ be a point on $\mathcal H$ excluding the arc $\overset{\frown}{AD}$ containing $B$.
Define $Q:=AC\cap PB$ and $R:=BD\cap PC$.

Then the area of the quadrilateral $BCRQ$ is independent of the position of $P$.
\end{proposition}

\begin{proof}
Since areas are scaled only by a constant factor under affine transformations,
we may assume $\kappa=1$, so that $\mathcal H:xy=1$.

We parametrize points on the hyperbola by
\[
X(t):=\Bigl(t,\frac{1}{t}\Bigr)\qquad (t>0).
\]
Then
\[
A=X(a),\quad B=X(ar),\quad C=X(ar^2),\quad D=X(ar^3),\quad P=X(p)\ (p>0).
\]

\medskip
\noindent
\textbf{Lemma 1 (Equation of a chord).}
For $t_1\neq t_2$, the line through $X(t_1)$ and $X(t_2)$ is given by
\[
y=\frac{(t_1+t_2)-x}{t_1t_2}.
\]
\emph{Proof.} This follows immediately from the two-point form. \qed

\medskip
Thus each chord is described by the sum $S=t_1+t_2$ and product $T=t_1t_2$:
\[
L(t_1,t_2):\quad y=\frac{S-x}{T}.
\]

\medskip
\noindent
\textbf{Lemma 2 (Intersection of two chords).}
Let
\[
L(t_1,t_2)\cap L(t_3,t_4)=:Z.
\]
Then
\[
x_Z=\frac{(t_3+t_4)\,t_1t_2-(t_1+t_2)\,t_3t_4}{t_1t_2-t_3t_4}.
\]
\emph{Proof.} Solve $\frac{S_{12}-x}{T_{12}}=\frac{S_{34}-x}{T_{34}}$. \qed

\medskip
Using this, we compute $Q$ and $R$.

For
\[
Q=L(a,ar^2)\cap L(p,ar),
\]
we obtain
\[
x_Q=\frac{a^2r^2+arp-ar^2p-ap}{ar-p},
\]
hence
\[
x_C-x_Q=\frac{a(ar^2-p)(r-1)}{ar-p}.
\]

Similarly, for
\[
R=L(ar,ar^3)\cap L(p,ar^2),
\]
we obtain
\[
x_R=\frac{a^2r^4+ar^2p-ar^3p-arp}{ar^2-p},
\]
hence
\[
x_B-x_R=\frac{a r^3(ar-p)(r-1)}{ar^2-p}.
\]

\medskip
We now compute the area structurally.
Since $Q,C$ lie on $AC$ and $R,B$ lie on $BD$,
we have
\[
[\,BCRQ\,]=\frac{1}{2}\,|\overrightarrow{QC}\times\overrightarrow{RB}|.
\]

Taking direction vectors
\[
\vec d_{AC}=(1,-\tfrac{1}{a^2r^2}),\qquad
\vec d_{BD}=(1,-\tfrac{1}{a^2r^4}),
\]
we obtain
\[
\overrightarrow{QC}=(x_C-x_Q)\vec d_{AC},\qquad
\overrightarrow{RB}=(x_B-x_R)\vec d_{BD}.
\]

Thus
\[
[\,BCRQ\,]
=
\frac{1}{2}\,|x_C-x_Q|\,|x_B-x_R|\cdot|\vec d_{AC}\times\vec d_{BD}|.
\]

A direct computation gives
\[
\vec d_{AC}\times\vec d_{BD}
=
\frac{r^2-1}{a^2r^4}.
\]

On the other hand,
\[
|x_C-x_Q|\,|x_B-x_R|
=
a^2r^3(r-1)^2,
\]
where the dependence on $p$ cancels.

Therefore,
\[
[\,BCRQ\,]
=
\frac12\cdot a^2r^3(r-1)^2\cdot \frac{|r^2-1|}{a^2r^4}
=
\frac{(r+1)|r-1|^3}{2r^2},
\]
which is independent of $p$.
\end{proof}

\end{document}
